% Two-page paper. Compile with the official CVPR author kit
% (https://github.com/cvpr-org/author-kit): put this file, refs.bib and figures/
% next to cvpr.sty and ieeenat_fullname.bst.
%
% Every \todo{...} marks a number or sentence that depends on the final
% experiments. Nothing marked \todo may remain in the submitted PDF.
\documentclass[10pt,twocolumn,letterpaper]{article}
\usepackage[pagenumbers]{cvpr}

\usepackage{amsmath,amssymb}
\usepackage{booktabs}
\usepackage{graphicx}
\usepackage{xcolor}
\usepackage[pagebackref,breaklinks,colorlinks,allcolors=cvprblue]{hyperref}

\newcommand{\todo}[1]{\textcolor{red}{[#1]}}
\newcommand{\Lcf}{\mathcal{L}_{\mathrm{cf}}}
\newcommand{\Lgrasp}{\mathcal{L}_{\mathrm{grasp}}}

\def\paperID{****}
\def\confName{FFI}
\def\confYear{2026}

\title{Does the Model Listen? Counterfactual Prompts for Language-driven Grasp Detection}

\author{\todo{First Last}\\
\todo{Affiliation}\\
{\tt\small \todo{email}}\\
{\small Code: \url{https://github.com/NhatBarbie20/QACI-Language-driven-Grasp-Detection}}
}

\begin{document}
\maketitle

\begin{abstract}
Language-driven grasp detection asks a model to grasp the object named in an
instruction. We observe that the usual success rate does not show whether the
instruction was used at all: a model can be right by grasping the most salient
object. We make two contributions. First, a diagnostic protocol that queries
each test image a second time with an instruction for a different object and
measures how often the prediction fails to move. Second, a counterfactual prompt
loss that uses the same swap during training and penalises a network whose
grasp quality map does not depend on the instruction. On Grasp-Anything++ a
lightweight FiLM-conditioned network trained with this loss reduces the ignore
rate from \todo{xx.x}\% to \todo{xx.x}\% \todo{while keeping / and improves} the
success rate. \todo{Rewrite the last sentence once results are final.}
\end{abstract}

\section{Introduction}
Grasp detection predicts where a parallel gripper should close on an
object~\cite{jiang2011rectangle,morrison2018ggcnn,kumra2020grconvnet}.
In cluttered scenes the robot must also know \emph{which} object is meant, and
language is the natural way to say so. Grasp-Anything++~\cite{vuong2024lgd}
made this setting trainable at scale by pairing images with grasp instructions
and rectangle labels.

A language-conditioned detector has an easy shortcut. Most images contain one
dominant object, so grasping it is often scored as correct whatever the
instruction says. The standard metric counts a prediction as a success if it
overlaps a ground-truth rectangle of the queried object, and therefore cannot
separate a model that follows the instruction from one that ignores it.

We address this in evaluation and in training.
\textbf{(i)}~We propose a prompt-swap protocol with two metrics, \emph{ignore
rate} and \emph{paired success rate}, that quantify language use directly.
\textbf{(ii)}~We propose a counterfactual prompt loss that contrasts, on the
same image, the instruction of the target object with the instruction of
another object, at no extra annotation cost.
\textbf{(iii)}~We report a controlled comparison of fusion strategies in a
2M-parameter network, with code and a reproducible data pipeline.

\section{Related Work}
\noindent\textbf{Grasp detection.}
Generative approaches such as GG-CNN~\cite{morrison2018ggcnn} and
GR-ConvNet~\cite{kumra2020grconvnet} predict dense quality, angle and width
maps, and remain strong and fast baselines. Grasp-Anything~\cite{vuong2024graspanything}
scaled the training data with foundation models.

\noindent\textbf{Language-driven grasping.}
LGD~\cite{vuong2024lgd} casts the task as conditional diffusion with a
contrastive objective defined on the denoising process, and later work adds
mask-guided attention~\cite{vo2024maskgrasp}. CLIPort~\cite{shridhar2021cliport}
fuses CLIP~\cite{radford2021clip} features for manipulation policies. Closest to
us, Nguyen~\etal~\cite{nguyen2024negativeprompt} learn a negative prompt
embedding that steers a 6-DoF diffusion model away from unwanted objects. We
differ in three ways: our negative is an explicit instruction of another object
in the same image, it acts through a margin on a dense 2D quality map and needs
no generative model, and we turn the same swap into an evaluation metric.

\begin{figure*}[t]
  \centering
  \includegraphics[width=0.92\linewidth]{figures/method.pdf}
  \caption{\textbf{Overview.} A convolutional stem encodes the image once. The
  residual blocks are modulated by a frozen CLIP embedding of the instruction
  through FiLM. During training the blocks are run a second time with an
  instruction $t^{-}$ for another object of the same image (dashed), and
  $\Lcf$ requires the quality inside the target grasp region $R$ to be higher
  for the matching instruction $t^{+}$.}
  \label{fig:method}
\end{figure*}

\section{Method}
\noindent\textbf{Grasp representation.}
A grasp is a rectangle $g=(x,y,w,h,\theta)$. Following~\cite{morrison2018ggcnn,kumra2020grconvnet},
the network $f$ maps an image $I$ and instruction $t$ to four maps: a quality
map $Q\in[0,1]^{H\times W}$, $\cos2\theta$, $\sin2\theta$ and a width map $W$.
The grasp is read at $(x^{*},y^{*})=\arg\max Q$ with
$\theta=\tfrac12\operatorname{atan2}(\sin2\theta,\cos2\theta)$.

\noindent\textbf{Language conditioning.}
The instruction is encoded by the frozen CLIP text encoder into
$\mathbf{e}_t\in\mathbb{R}^{512}$. Each of the five residual blocks of a
GR-ConvNet-style backbone is modulated with FiLM~\cite{perez2018film}:
\begin{equation}
  \mathbf{F}' = \bigl(1+\gamma(\mathbf{e}_t)\bigr)\odot\mathbf{F} + \beta(\mathbf{e}_t),
  \label{eq:film}
\end{equation}
where $\gamma,\beta$ are linear maps initialised to zero, so training starts
from the vision-only network. The stem before the first block does not depend
on $t$ and is shared between instructions (Fig.~\ref{fig:method}).

\noindent\textbf{Counterfactual prompt loss.}
For an image with target instruction $t^{+}$ we sample an instruction $t^{-}$
of a different annotated object in the same image. Let $R$ be the ground-truth
grasp region of $t^{+}$ and $\bar{Q}_R(t)=\tfrac{1}{|R|}\sum_{p\in R}Q_p(I,t)$. We require
\begin{equation}
  \Lcf = \max\bigl(0,\; m-\bar{Q}_R(t^{+})+\bar{Q}_R(t^{-})\bigr).
  \label{eq:cf}
\end{equation}
Since both terms share the image, only the language pathway can reduce
$\Lcf$; a network that ignores $t$ always pays the full margin $m$. The total
objective is
\begin{equation}
  \mathcal{L} = \Lgrasp + \lambda\,\Lcf ,
\end{equation}
where $\Lgrasp$ is a binary cross-entropy on $Q$ plus smooth-$\ell_1$ losses on
the other three maps. We use $m=0.5$ and $\lambda=1$. The second pass reuses the
stem features, which keeps the overhead to the residual blocks and decoder.

\noindent\textbf{Prompt-swap metrics.}
At test time we predict $g^{+}=f(I,t^{+})$ and $g^{-}=f(I,t^{-})$. The
\emph{ignore rate} is the fraction of correctly solved samples for which
$g^{-}$ still matches the ground truth of $t^{+}$; it is $1$ for any
vision-only model. The \emph{paired success rate} (PSR) is the fraction of
samples where $g^{+}$ is correct and $g^{-}$ is not on the first object.

\section{Experiments}
\noindent\textbf{Setup.}
We use the official LGD splits of Grasp-Anything++~\cite{vuong2024lgd}. After
removing objects without a language annotation in the public release and 37
training scenes that also occur in the unseen test list, this gives 26{,}820
training, 1{,}363 seen and 597 unseen test samples; the 25 unseen categories
never appear in training. Images are resized to $224^2$. All models are trained
for \todo{30} epochs with AdamW on one T4 GPU. We report the success rate
(IoU${>}25\%$, angle error ${<}30^{\circ}$)~\cite{jiang2011rectangle}, its
harmonic mean $H$ over seen and unseen, and the two prompt-swap metrics, which
are computed on the \todo{xx}\% of test images that contain a second annotated object.

\begin{table}[t]
  \centering
  \small
  \setlength{\tabcolsep}{3.5pt}
  \begin{tabular}{@{}lccccc@{}}
    \toprule
    & \multicolumn{3}{c}{Success rate $\uparrow$} & Ignore $\downarrow$ & PSR $\uparrow$ \\
    \cmidrule(lr){2-4}
    Method & Seen & Unseen & $H$ & Unseen & Unseen \\
    \midrule
    \multicolumn{6}{@{}l}{\emph{Reported in~\cite{vuong2024lgd} (different test subset)}}\\
    GR-ConvNet + CLIP & 0.37 & 0.18 & 0.24 & -- & -- \\
    LGD (diffusion)   & 0.48 & 0.42 & 0.45 & -- & -- \\
    \midrule
    \multicolumn{6}{@{}l}{\emph{Trained by us, same data and protocol}}\\
    Vision only       & \todo{.xx} & \todo{.xx} & \todo{.xx} & 1.00 & 0.00 \\
    Concat fusion     & \todo{.xx} & \todo{.xx} & \todo{.xx} & \todo{.xx} & \todo{.xx} \\
    FiLM fusion       & \todo{.xx} & \todo{.xx} & \todo{.xx} & \todo{.xx} & \todo{.xx} \\
    FiLM + $\Lcf$ (ours) & \todo{.xx} & \todo{.xx} & \todo{.xx} & \todo{.xx} & \todo{.xx} \\
    \bottomrule
  \end{tabular}
  \caption{Results on Grasp-Anything++. The top rows are copied from the
  original paper and are not directly comparable, because part of the official
  split has no public language annotation.}
  \label{tab:results}
\end{table}

\noindent\textbf{Results.}
\todo{Three or four sentences, written after the runs finish. State (1) what the
vision-only row shows about the shortcut, (2) the effect of FiLM over
concatenation, (3) the effect of $\Lcf$ on ignore rate and PSR, and on the
success rate, including the case where it does not help. Add one qualitative
figure (same image, two instructions) if space allows.}

\noindent\textbf{Limitations.}
In 68\% of multi-part objects of our subset every part carries identical grasp
labels, so the data supports object-level but not part-level grounding. The
prompt-swap metrics assume that the grasps of two objects do not overlap.
\todo{Add: single seed / small backbone, as applicable.}

\section{Conclusion}
Success rate alone hides whether a language-driven grasp detector uses its
instruction. Swapping the prompt on a fixed image exposes this at test time and
supplies a training signal for free. \todo{One sentence with the main number.}
Future work includes part-level negatives and applying the loss to
diffusion-based detectors.

{\small
\bibliographystyle{ieeenat_fullname}
\bibliography{refs}
}
\end{document}
